\documentclass{article}
\usepackage{amsmath,amssymb,amsthm,mathtools}
\newtheorem{proposition}{Proposition}
\newcommand{\CC}{\mathcal{C}}

% A test of notions defined in the running text (docs/writing.md): `\(…\)` math,
% a notion whose condition is not read (its subject still counts), a notion by a
% monotone map without "the map", and four ways of using a notion in a result.
% The expected output of `proofread check --vocab unitsquare --claims` is check.txt.

\begin{document}

\section{Notions}

Let $\CC$ denote the class of bivariate copulas.
We call a copula \(C\) \emph{stochastically increasing} (\emph{SI}) if its Markov kernel has a version such that, for every \(v\in(0,1)\), the map \(t\mapsto K_C(t,[0,v])\) is non-increasing. We also use a classical tail notion. \(C\) is \emph{left-tail decreasing} (\emph{LTD}) if, for every \(v\in(0,1)\),
\begin{equation}\label{eq:ltd}
u\longmapsto\frac{C(u,v)}{u},\qquad u\in(0,1],
\end{equation}
is non-increasing.

\section{Using the notion}

\begin{proposition}
For every LTD copula $C$, $C(u,v)\ge uv$ for all $u,v\in[0,1]$.
\end{proposition}

\begin{proposition}
Let $C$ be an LTD copula. Then $C(u,v)\ge uv$ for all $u,v\in[0,1]$.
\end{proposition}

\begin{proposition}
If $C$ is an LTD copula, then $C(u,v)\ge uv$ for all $u,v\in[0,1]$.
\end{proposition}

\begin{proposition}
Every LTD copula $C$ satisfies $C(u,v)\ge uv$ for all $u,v\in[0,1]$.
\end{proposition}

\end{document}
